\chapter{Conclusion}
\label{chap:conclusion}

% Template

% You generally cover three things in the Conclusions section, and each of these usually merits a separate subsection:

% \begin{enumerate}
% 	\item Conclusions
% 	\item Summary of Contributions
% 	\item Future Research
% \end{enumerate}

% Conclusions are not a rambling summary of the thesis: they are short, concise statements of the inferences that you have made because of your work. It helps to organize these as short numbered paragraphs, ordered from most to least important. All conclusions should be directly related to the research question stated in the introduction. Examples:

% \begin{enumerate}
% 	\item The problem stated in the introduction has been solved: as shown in Sections ? to ??, an algorithm capable of handling large-scale Zylon problems in reasonable time has been developed.
% 	\item The principal mechanism needed in the improved Zylon algorithm is the Grooty mechanism.
% 	\item Etc.
% \end{enumerate}

% The Summary of Contributions will be much sought and carefully read by the examiners. Here you list the contributions of new knowledge that your thesis makes. Of course, the thesis itself must substantiate any claims made here. There is often some overlap with the Conclusions, but that's okay. Concise numbered paragraphs are again best. Organize from most to least important. Examples:

% \begin{enumerate}
% 	\item Developed a much quicker algorithm for large-scale Zylon problems.
% 	\item Demonstrated the first use of the Grooty mechanism for Zylon calculations.
% 	\item Etc.
% \end{enumerate}

% The Future Research subsection is included so that researchers picking up this work in future have the benefit of the ideas that you generated while you were working on the project. Again, concise numbered paragraphs are usually best. 

\section{Conclusions}
\label{sec:conclusions}

In this thesis we have shown that simply enforcing low-dimensional structures or latent layers does not
improve the performance of controllers for overparameterized systems; instead, it might even hurt performance. However,
when combining sufficiently pretrained latent projectors with the right controller architecture, they can
achieve similar performances to fully connected networks, while being more efficient in their parameter count
and training time. While this could be clearly shown for LQR, under RL only tendencies were observed, requiring
further investigation.

We further found conflicting results on whether low-dimensional latent states
can improve robustness to measurement noise. While for LQR solvers and model-informed
\emph{StE} this holds true, for the camera environment, we could not clearly observe that. This however might
be related to the first point, where our RL approaches showed already bad performance
under the clean loop.

Regarding the latent state, it is also important to ensure that not only control information is carried
here, but also information about the observed state. Otherwise, robustness to control noise might worsen.

In relation to that, we also note that a low degradation in control noise can also hint at a barely acting
controller, and is not always a sign of robustness. Still, in our experiments, low-dimensional controllers 
were affected slightly more by it than fully parameterized ones, which we observed across all experiments.

\section{Summary of Contributions}
\label{sec:contributions}

We have presented that combining pretrained encoders with a trainable network
can reduce the parameter count effectively on a toy example, while noting common problems in relation to
architecture design. Failing to translate the success of it onto a vision-based problem brought
additional pitfalls to light. Further, we showed that fully parameterized controllers offer, if even only slightly, more
robustness to control noise, while for the other noise types we could only observe tendencies,
which open possibilities for further research here.

\section{Future Research}
\label{sec:future-research}

From where this thesis ends, there are different directions, which we have left for future research
and discussion. These include, but are not limited to, whether:

\begin{enumerate}
	\item and how the number of parameters can be distributed more effectively between
	the extractor and the actor, to achieve performances comparable to a fully connected network;
	\item our results, such as better robustness to measurement noise of low-dimensional LQR controllers, might hold
	for more complex systems, i.e. a walking humanoid robot;
	\item these findings are transferable to real-world applications, i.e., if this setup
	were built and tested;
\end{enumerate}

Finally, we also want to mention that, while these are interesting directions, the data produced from
our experiments still offers a lot of possibilities for additional analysis. For example,
while we captured the compute time and environment steps needed, a different angle could include
this in the evaluation of the approaches. And while it seems like our conclusions
leave us with more questions than we started with, this is what research is for.